\section{Algorithme de Broker Charles et Lauter}\label{SECTION friableEtBroker}

L'objectif de cette section est d'énoncer l'algorithme de Broker Charles et Lauter \cite{BrokerCharlesLauter}, optimisé par Jao et Soukharev \cite{JaoSoukharev}, et dont je propose une implémentation. Il prend en entrée une courbe elliptique ordinaire $E$ définie sur $\F_{q}$ de caractéristique $p$, son anneau d'endomorphismes $\OR_{\Delta}$, un point $P\in E(\F_{q^{n}})$ avec $n\geq1$ et un idéal $\LR$ de norme $l$ premier, et qui calcule la $l$-isogénie $\phi$ résultant de l'action de l'idéal $\LR$ sur $E$. L'entrée de l'algorithme doit vérifier l'hypothèse \ref{HypothèseInversion}, pour que l'isogénie $\phi$ soit horizontale, et pour que l'on puisse évaluer des endomorphismes sur des points de $E(\F_{q^{n}})$ (voir propositions \ref{evalEndoAvecDenom} et \ref{KohelStructureVolcan}). 

\begin{Hyp}\label{HypothèseInversion}
	Soit $f$ le conducteur de $E$, $f_m$ son conducteur maximal sur $\F_{q}$. On suppose que $E$, $l$ et $n$ sont choisis tel que $l$ ne divise pas $A = pf_m|E(\F_{q^{n}})|$, que $l$ est décomposé dans $End(E)$ et que $f_{m}$ et $|E(\F_{q^{n}})|$ sont premiers entre eux.
\end{Hyp}

 Les méthodes pour calculer explicitement une isogénie sont présentées dans la section \ref{SECTION isogénieDegrésPetits}. Elles dérivent toutes des formules de Vélu données dans la proposition \ref{FormuleVélu}, et sont de complexité exponentielle en $l$. Cependant, dans la section \ref{SECTION isogenieFriable}, on explique comment ramener le calcul de la $l$-isogénie $\phi$ à celui de plusieurs isogénies de petits degrés. 
 
 Dans le fichier Broker.ipynb, je propose une implémentation prenant en compte l'optimisation de \cite{JaoSoukharev} ainsi que des initiatives personnelles. Mon approche est expliquée dans la section \ref{SECTION CalculfactoClasse}. J'analyse la complexité obtenue pour la comparer avec les résultats de \cite{JaoSoukharev} dans la section \ref{SECTION Broker}.   

\subsection{Isogénies horizontales de petits degrés}
\label{SECTION isogénieDegrésPetits}

\subsubsection{Algorithme SquareVélu}\label{SECTION squareVélu}

Étant donné un modèle de courbe elliptique $E$ défini sur $\F_{q}$, les formules de Vélu \ref{FormuleVélu} permettent de construire une forme explicite d'une $l$-isogénie $\phi$ à partir de la connaissance de son noyau. Appliquer directement les formules demande cependant $O(l)$ opérations dans $\F_{q}$. On peut obtenir une complexité en $O(\sqrt{l})$ selon la méthode de \cite{Velusqrt}. Il faut cependant préciser comment représenter le noyau. En effet, $\ker(\phi)$ peut être donné sous deux formes : un générateur $K\in E(\F_{q^{r}})$ avec $r \geq 1$, ou un polynôme $F$ dont les racines sont les abscisses des points du noyau. 

\begin{Def}
	Soit $E$ un modèle de courbe elliptique et $\phi : E\rightarrow E'$ une $l$-isogénie, avec $l$ premier impair. On sépare $\ker(\phi)$ en une partition $\left\lbrace \OR \right\rbrace \cup R\cup R'$, où $R$ regroupe $(l-1)/2$ points finis distincts, et $R'$ leur inverse. Un polynôme $F$ décrit le noyau de $\phi$ sur $E$ si c'est un polynôme unitaire de degré $(l-1)/2$ dont les racines sont les abscisses des points finis de $R$. Un tel polynôme ne dépend pas de la partition choisie. 
\end{Def}

Pour calculer une $l$-isogénie, \cite{Velusqrt} propose une méthode pour évaluer les polynômes de $\F_{q}[X]$ intervenant dans les formules de Vélu en $O(\sqrt{l})$ opérations dans $\F_{q}$. Les résultats suivants en découlent.

\begin{Prop}
	Soit $E$ un modèle de courbe elliptique défini sur $\F_{q}$, et $\phi : E\rightarrow E'$ une $l$-isogénie, avec $l$ premier impair. Soit $K$ un générateur de $Ker(\phi)$, et $P\in E(\F_{q})$. Il est possible de calculer un modèle de $E'$ et $\phi(P)\in E'$ en $O(\sqrt{l})$ opérations dans $\F_{q}$, en utilisant l'algorithme de \cite{Velusqrt} que l'on appelera SquareVélu.
\end{Prop}

\begin{Rem}
	SquareVélu est une amélioration asymptotique de l'application directe des formules de Vélu. Cependant pour les petits nombres premiers, l'application directe reste plus rapide. Avec l’implémentation de Sage, il est conseillé d'utiliser SquareVélu uniquement si $l > 100$. 
\end{Rem}

On décrit maintenant comment adapter les formules de Vélu à l'utilisation de polynômes décrivant le noyau d'une isogénie.

\begin{Prop}[\cite{Kohel}(\S 2.4)]
	
	Soit $E : y^2 = x^3 + ax + b$ un modèle de courbe elliptique défini sur $\F_{q}$, et $\phi$ une $l$-isogénie, avec $l$ premier impair. Soit $F$ le polynôme décrivant le noyau de $\phi$ sur $E$. Alors on peut obtenir une forme explicite de $\phi$ avec
	\begin{equation*}
		\phi(x,y) = \left( \frac{G(X)}{F(X)^2}\;,\;\frac{H(X,Y)}{F(X)^3}\right), 
	\end{equation*}
	où $G$ et $H$ peuvent être obtenus explicitement à partir de $F$. Pour $P\in E(\F_{q})$, on peut alors calculer $\phi(P)$ en $O(\sqrt{l})$ opérations dans $\F_{q}$, via les méthodes de \cite{Velusqrt}.
\end{Prop}

\begin{Rem}
	Il existe aussi des formules adaptées au cas où $l = 2$.
\end{Rem}




\subsubsection{Générateurs de groupes de torsion d'idéaux}\label{SECTION generateurTorsion}


Dans cette partie, on suppose connue une courbe elliptique $E$ définie sur $\F_{q}$ ordinaire avec $End(E)\simeq\OR_{\Delta}$ normalisé, et un idéal $\LR\subset\OR_{\Delta}$ de norme $l$ premier décomposé. On cherche un générateur du groupe de torsion $E\left[ \LR\right]\subset E(\Fb_{p})$. On suppose \ref{HypothèseInversion} donc $l$ est premier avec le conducteur maximal $f_{m}$ de $E(\F_{q})$.

\begin{Rem}
	 $\LR$ est un idéal primitif car de norme $l$ premier. La proposition \ref{formeIdeauxBaseFq} permet alors d'obtenir
	\begin{equation*}
		\LR = \left[ l, \frac{c + f\pi_{q}}{f_{m}}\right], \; c\in\Z.
	\end{equation*}
	Dans la suite on suppose $\LR$ donné sous cette forme. 
\end{Rem}


\begin{Prop}
	Avec les notations précédentes, $E[\LR]$ est un espace propre de l'endomorphisme $\pi_q$ restreint à $E[l]$, de valeur propre $cf^{-1}$ modulo $l$. 
\end{Prop}

\begin{proof}
	$E\left[ \LR\right]$ est par définition l'intersection de $E\left[ l\right] $ avec le noyau de $\left[\frac{c + f\pi_{q}}{f_{m}}\right]$. Or comme $l$ est premier avec $f_m$, l'endomorphisme $[f_m]$ restreint à $E[l]$ est inversible. Donc les restrictions à $E\left[ l\right]$ des endomorphismes $ \left[ \frac{c + f\pi_{q}}{f_{m}}\right]$ et $\left[ c + f\pi_{q}\right]$ ont même noyau. Si $P\in E[\LR]$, alors en particulier $f\pi_{q}(P) + cP = 0$, donc $\pi_{q}(P) = -cf^{-1}P$.
\end{proof}

On peut relier cette proposition avec la structure de volcan d'isogénie.

\begin{Prop}[\cite{MMSTV}(Prop 1)]
	
	Soit $E$ définie sur $\F_{q}$ de $j$-invariant différent de $0$, $1728$. Soit $l$ un nombre premier, différent de la caractéristique $p$. 
	\begin{enumerate}
		\item $E$ admet une unique $l$-isogénie définie sur $\F_{q}$ si et seulement si $\pi_{q}$ restreint à $E[l]$ a un unique espace propre de dimension $1$.
		\item $E$ admet exactement deux $l$-isogénies définies sur $\F_{q}$ si et seulement si $\pi_{q}$ restreint à $E[l]$ a deux valeurs propres distinctes.
		\item $E$ admet $l+1$ isogénies définies sur sur $\F_{q}$ si et seulement si $\pi_{q}$ restreint à $E[l]$ a une valeur propre double.
	\end{enumerate}
	
\end{Prop}

Puisque $l$ est supposé décomposé et premier avec $f_{m}$, on sait qu'il existe exactement deux $l$-isogénies définies sur $\F_{q}$. On note $v_{1}$ et $v_{2}$ les deux valeurs propres de $\pi_{q}$ restreint à $E[l]$. 

\begin{Prop}[\cite{MMSTV}(Prop 2)]
	Soit $G_{1}$, $G_{2}$ les deux sous-groupes de $E$, noyaux des deux $l$-isogénies définies sur $\F_{q}$. Soit $r_{1}$, $r_{2}$ minimaux tels que $G_{i}\subset E\left( \F_{q^{r_{i}}}\right) $. Alors $r_{i}$ est l'ordre de $v_{i}$ dans $\F_{l}^{*}$ (en particulier $r_{i}$ divise $l-1$ ). De plus si $r_{i}$ est pair, les abscisses des points de $G_{i}$ sont dans $\F_{q^{r_{i}/2}}$.
\end{Prop}

La proposition précédente nous donne alors le degré de l'extension dans laquelle on doit chercher $E[l]$ en résolvant un problème de logarithme discret dans $\F_{l}^{*}$, que l'on note $r$. L'arithmétique dans l'extension $\F_{q^{r}}$ est de complexité en $O(r\log q)$, donc exponentielle en $r$. Or il est possible que $r$ soit de taille de l'ordre de $l$. Par la suite, on voudra dès que possible mener des calculs uniquement sur le corps de base $\F_{q}$. On évitera donc de représenter le noyau $E[\LR]$ par un générateur, mais plutôt par un polynôme.


\subsubsection{Représentation polynomiale de groupes de torsion d'idéaux.}\label{SECTION méthodeElkies}

 On reprend les notations de la section précédente.
 On présente ici une méthode pour représenter un groupe de torsion d'un idéal $\LR$ par un polynôme. Elle a l'avantage d'éviter tout calcul dans des extensions de $\F_{q}$. En contrepartie, cette méthode impose des conditions sur la taille du degré $l$. On résume l'idée de la méthode en admettant les résultats utilisés.
\begin{enumerate}
	
	\item Il existe un polynôme $\Psi_{l}(X,Y)$, appelé polynôme modulaire, dépendant uniquement de $l$, tel que les racines de $\Psi_{l}(j,Y)$ dans $\Fb_{p}$ sont exactement les $j$-invariants des courbes $l$-isogènes à $E$. Les racines appartenant à $\F_{q}$ correspondent aux courbes images des $l$-isogénies définies sur $\F_{q}$.
	
	\item Soit $h\in\F_{q}$ une racine de $\Psi_{l}(j,Y)$, alors il est possible d'obtenir un modèle défini sur $\F_{q}$ de la courbe elliptique $E'$ de $j$-invariant $h$, à partir de la donnée de $E$, $h$ et $l$, en un nombre constant d'opérations dans $\F_{q}$. On utilise des formules explicites qui imposent
	\begin{equation*}
		l < 4|\Delta|.
	\end{equation*}
	
	\item Soit $F_{h}$ le polynôme qui décrit le noyau de la $l$-isogénie $\phi: E\rightarrow E'$. Soit $\tilde{F}_{h}$ défini de la même façon à partir d'un relèvement de Deuring de $E$, $E'$ et $\phi$. On peut alors retrouver explicitement les coefficients de $\tilde{F}_{h}$. Si chaque terme a une bonne réduction modulo $p$, alors on obtient les coefficient de $F_{h}$. La bonne réduction n'est cependant pas garantie, il faut que $l$ vérifie
	\begin{equation*}
		p > 4(l+1).
	\end{equation*}
	Cette méthode a été proposée par Elkies, et amélioré dans \cite{Bostan_2008}. Les améliorations visent principalement à minimiser et optimiser des calculs de séries de Laurent nécessaires pour obtenir $\tilde{F}_{h}$. En suivant la convention utilisée par Sage, on appellera BMSS l'algorithme proposé dans \cite{Bostan_2008}(\S 4.3, fastElkies') calculant $F_{h}$ à partir de $E$, $E'$ et $l$.


	\item D'après l'hypothèse $\ref{HypothèseInversion}$, $\Psi_{l}(j,Y)$ admet exactement deux racines dans $\F_{q}$. On souhaite trouver l'unique isogénie normalisée de noyau $E[\LR]$. Il faut donc tester si le polynôme $F$ obtenu par BMSS décrit le bon noyau. Pour cela, on vérifie formellement la valeur propre avec laquelle agit le morphisme de Frobenius sur l'ensemble de points décrit par $F$, comme énoncé dans la proposition \ref{propositionCheckElkies}.

\end{enumerate}

\begin{Prop}\label{propositionCheckElkies}
	Soit $E\in Ell_{\F_{q}}(\OR_{\Delta})$. Soit $\LR = \left[ l, \frac{c + f\pi_{q}}{f_{m}}\right]$ un idéal de $\OR_{\Delta}$ de norme $l$ premier avec le conducteur maximal $f_m$, et $k$ l'inverse de $f$ modulo $l$. Soit $F\in\F_{q}[X]$. On pose 
	\begin{equation*}
		R = \F_{q}\left[ X,Y\right] / \left( F, Y^2 - X^3 - aX - b\right).
	\end{equation*}
	Soit $(R_{1}(X), YR_{2}(X))$ les fonctions rationnelles qui définissent l'endomorphisme $[-ck]$. Alors $F$ est le polynôme décrivant $E\left[ \LR\right]$ si et seulement si, dans le corps de fractions de $R$, on a:
	\begin{equation*}
		\bar{X}^q = R_{1}(\bar{X})\;,\;\bar{Y}^{q} = \bar{Y}R_{2}(\bar{X}).
	\end{equation*}
\end{Prop}

\begin{Rem}
	Dans toute la suite, pour vérifier si un polynôme donné décrit le bon noyau, on utilise l'algorithme CheckElkies de l'implémentation Sage de \cite{Pegasis}. J'ai aussi essayé d'implémenter cet algorithme, mais ma version n'est pas aussi bien optimisée.  
\end{Rem}

\begin{Hyp}\label{HypothèseElkies}
	Soit $E\in Ell_{\F_{q}}(\OR_{\Delta})$. Un nombre premier $l$ vérifie l'hypothèse \ref{HypothèseElkies} si $l < 4|\Delta|$ et $p > 4(l+1)$.
\end{Hyp}

On en déduit l'algorithme \ref{ALGO BMSSCheckElkies} permettant de calculer le polynôme décrivant le sous-groupe de $\LR$ torsion d'un courbe $E\in Ell_{\F_{q}}(\OR_{\Delta})$. On suppose ici que le polynôme modulaire $\Psi_{l}$ est pré-calculé.  

\begin{algorithm}[h]
	\caption{Méthode d'Elkies}\label{ALGO BMSSCheckElkies}
	\begin{algorithmic}[1]
		\Require Une courbe elliptique $E\in Ell_{\F_{q}}(\OR_{\Delta})$, de conducteur maximal $f_m$. Un nombre $l$ premier  vérifiant les hypothèses \ref{HypothèseInversion} et \ref{HypothèseElkies}. Un idéal $\LR = \left[ l, \frac{c + f\pi_{q}}{f_{m}}\right]$. Le polynôme modulaire $\Psi_{l}$.
		\Ensure Un polynôme $F\in\F_{q}\left[ X\right]$ décrivant $E[\LR]$. 
		\State $P(Y)\gets PGCD(\Psi_{l}(j,Y),Y^q - Y)$ ($P$ est de degré 2 scindé.)
		\State $h_{1}, h_{2}\gets$ racines de $P$
		\State $E_{1}\gets$ modèle normalisé de $h_{1}$ (formules explicites selon $E$, $l$ et $h_1$)
		\State $F_{1}\gets BMSS(E,E_{1},l)$
		\If{$F_{1}$ décrit $E[\LR]$} 
		\State \Return $F_{1}$
		\Else 
		\State $E_{2}\gets$ modèle normalisé de $h_{2}$ (formules explicites selon $E$, $l$ et $h_1$)
		\State $F_{2}\gets BMSS(E,E_{2},l)$
		\State \Return $F_{2}$
		\EndIf
	\end{algorithmic}
\end{algorithm}

\begin{Rem}
	Lorsque l'on calcule à la suite plusieurs $l$-isogénies horizontales associées à l'action d'un même idéal $\LR$, on peut éviter d'avoir à deviner le bon $j$-invariant avec CheckElkies en retenant les $j$-invariants vus aux étapes précédentes.  
\end{Rem}

\subsection{Calculs d'isogénies par équivalence friable}\label{SECTION isogenieFriable}
%Ici problème de l'idéal principale, explication des conditions sur le conducteur/cardinal
%Calculs de composantes, sans analyses de complexité générale (lemmes ?)

Avec ce qui précède, on sait uniquement calculer des $l$-isogénies avec $l$ vérifiant les hypothèses \ref{HypothèseInversion} et \ref{HypothèseElkies}, avec une complexité exponentielle en $l$. 

\begin{Def}
	Soit $E\in Ell_{\F_{q}}(\OR_{\Delta})$. Un idéal $\LR\in\OR_{\Delta}$ est dit friable s'il vérifie 
	\begin{equation*}
		\LR = \PR_{p_{1}}^{e_{1}}\PR_{p_{2}}^{e_{2}}\ldots\PR_{p_{N}}^{e_{N}},
	\end{equation*} 
	où $e_{i}\in\N$ et les idéaux $\PR_{p_{i}}\in\OR_{\Delta}$ sont de norme $p_{i}$, appartenant à une liste $L_{premier}$ de nombres premiers vérifiant les hypothèses \ref{HypothèseInversion} et \ref{HypothèseElkies}.
\end{Def}

On se donne un idéal $\LR\in\OR_{\Delta}$, et on suppose que l'on peut trouver un idéal entier $\LR'$ équivalent à $\LR$ et friable. On peut donc calculer l'action de $\LR'$ sur $E$ en calculant $e_1 + \ldots + e_N$ actions d'idéaux avec l'algorithme \ref{ALGO BMSSCheckElkies}. On souhaite en déduire l'action de $\LR$ sur $E$. 

\begin{Def}
	 Avec les notations précédentes, soit $\phi_\LR : E \rightarrow E/E[\LR]$. On appelle composantes de $\phi_\LR$ associées à $\LR'$ les isogénies obtenues par la suite
	\begin{equation*}
		E\rightarrow E/E\left[ \PR_{1}\right] \rightarrow E/E\left[ \PR_{1}^{2}\right] \rightarrow\ldots\rightarrow E/E\left[ \PR_{1}^{e_{1}}\right] \rightarrow
	\end{equation*}
	\begin{equation*}
		\rightarrow E/E\left[ \PR_{1}^{e_{1}}\PR_{2}\right] \rightarrow\ldots\rightarrow E/E\left[ \PR_{1}^{e_{1}}\PR_{2}^{e_{2}}\right] \rightarrow\ldots\rightarrow E/E\left[\PR_{1}^{e_{1}} \ldots\PR_{N}^{e_{N}}\right].
	\end{equation*}
	On note $\phi_c : E\rightarrow E_c$ la composée de toutes les composantes.
\end{Def}

$E_c$ est donc la courbe elliptique $E/E[\LR]$. Cependant la $l$-isogénie $\phi_\LR$ n'est pas égale à $\phi_c$, qui n'est pas de degré $l$ a priori. On a cependant la relation
\begin{equation*}
	\LR = (\alpha\OR_{\Delta})\PR_{p_{1}}^{e_{1}}\PR_{p_{2}}^{e_{2}}\ldots\PR_{p_{N}}^{e_{N}},
\end{equation*}
où $\alpha\OR_\Delta$ est un idéal fractionnaire principal. A priori $\alpha\in\K$, cependant
\begin{equation*}
	\LR\bar{\PR_{1}}^{e_{1}}\bar{\PR_{2}}^{e_{2}}\ldots\bar{\PR_{N}}^{e_{N}} = \beta\OR_{\Delta},
\end{equation*}
où $\beta\OR_{\Delta}$ est un idéal principal entier. Ainsi $\beta\in\OR_{\Delta}$ et $\alpha = \frac{\beta}{m}$ avec $m = p_1^{e_1}\ldots p_N^{e_N}$. On calcule $\beta$ en utilisant l'algorithme \ref{ALGO idealprincipal}, et on trouve ainsi $\alpha = \frac{u + v\pi_{q}}{mf_{m}}\in\OR_{\Delta}$, $u$, $v\in\Z$. D'après l'hypothèse \ref{HypothèseInversion}, les nombres premiers $p_{i}$ sont premiers avec $|E(\F_{q^{n}})|$, donc $m$ et $f_{m}$ le sont aussi, ce qui permet de calculer $k = (mf_{m})^{-1}$ modulo $|E(\F_{q^{n}})|$. L'endomorphisme correspondant à l'action de l'idéal $\alpha\OR_{\Delta}$ est donc $k(u + v\pi_{q})$. D'après la proposition \ref{facteurEndomorphismeReduction}, son facteur est $ku$ modulo $p$. On en déduit l'algorithme \ref{ALGO IsogenieEquivalenceFriable} qui calcule l'action d'un idéal à partir de celle d'un idéal équivalent friable, sous forme explicite normalisée. 

\begin{algorithm}[H]
	\caption{Calcul d'isogénies par équivalence friable}\label{ALGO IsogenieEquivalenceFriable}
	\begin{algorithmic}[1]
		\Require $E\in Ell_{\F_{q}}(\OR_{\Delta})$. un entier positif $n$ et un nombre premier $l$ vérifiant l'hypothèse \ref{HypothèseInversion}. Un idéal $\LR\subset\OR_{\Delta}$ de norme $l$. Un point $P\in E(\F_{q^{n}})$. Un idéal $\LR'$ friable équivalent à $\LR$, $L_{facteurs}$ la liste des idéaux premiers de la factorisation de $\LR'$, et $L_{exposants}$ la liste des exposants associés.
		\Ensure Un modèle normalisé $E'$ de la courbe elliptique image de la $l$-isogénie $\phi_\LR$. L'image $\phi_\LR(P)$ du point $P$. 
		\State $L_{composantes}\gets$\O
		\State $E_{d}\gets E$
		\State $\LR_{principal}\gets\LR$
		\State $Q\gets P$
		\For {$e_{i}$, $\PR_{i}$ parcourant $L_{exposants}$, $L_{facteurs}$}
		\State$\LR_{principal}\gets\LR_{principal}\times\bar{\PR_{i}}^{ei}$
		\State$\Psi\gets \Psi_{p_{i}}$ le polynôme modulaire de $p_{i}$ norme de $\PR_{i}$. 
		\For {$j\in\left[ 1\ldots e_{i}\right]$ }
		\State Calculer $\phi_c : E_{d} \rightarrow E_{d}/ E_{d}[\PR_{i}]$ avec l'algorithme \ref{ALGO BMSSCheckElkies}.
		\State $E_c\gets$ modèle explicite normalisé de $E_{d}/ E_{d}[\PR_{i}]$ 
		\State $E_{d}\gets E_c$ 
		\State $Q\gets \phi_c(Q)$
		\State Mettre $\phi_c$ dans $L_{composantes}$
		\EndFor
		\EndFor
		\State $\beta\gets$ générateur de l'idéal principal $\LR_{principal}$ via l'algorithme \ref{ALGO idealprincipal}.
		\State Calculer $u$, $v$ tels que $\beta = \frac{u + v\pi_{q}}{f_{m}}$ (proposition \ref{formeIdeauxBaseFq}).
		\State $m\gets p_{1}^{e_{1}}\ldots p_{g}^{e_{g}}$, où $\PR_{i}$ est de norme $p_{i}$ 
		\State $k\gets (mf_m)^{-1}$ modulo $|E(\F_{q^{n}}|)$
		\State $c\gets uk$
		\State $a$, $b\gets$ coefficient du modèle $E_c : y^2 = x^3 + ax + b$
		\State Poser $E' : y^2 = x^3 + c^4ax + c^6b$
		\State $Q\gets (k(u + v\pi_{q}))(Q)$
		\State \Return $E'$, $Q$
	\end{algorithmic}
\end{algorithm}

\begin{Prop}[\cite{Bostan_2008}]
	
	Soit $M : \N\times\N\rightarrow\N$ une fonction telle que $M(n,q)$ soit la complexité de la multiplication de deux polynômes de degrés au plus $n$ dans $\F_{q}[X]$. L'algorithme BMSS calcule le polynôme décrivant le noyau d'une isogénie définie sur $\F_{q}$ de degré $p_{i}$ avec une complexité en 
	\begin{equation*}
		O\left( M(p_{i},q)\log p_{i}\right).
	\end{equation*}
\end{Prop}

\begin{Prop}
	La méthode d'Elkies (algorithme \ref{ALGO BMSSCheckElkies}), appliquée à une courbe elliptique $E$ et un idéal $\PR_{i}\subset\OR_{\Delta}$ de norme $p_{i}$, retourne un polynôme $F$ décrivant $E[\PR_{i}]$ avec une complexité en 
	\begin{equation*}
		O\left( log(q) + M(p_{i},q)\log p_{i}\right).
	\end{equation*}
\end{Prop}

\begin{proof}
	L'algorithme d'Euclide pour calculer le PGCD à la ligne $1$ se fait en $O(\log q)$. L'étape suivante est une application de l'algorithme BMSS. Enfin, il faut vérifier le choix de $j$-invariant avec la proposition \ref{propositionCheckElkies}.
	La complexité de cette étape est dominée par les divisions de polynômes nécessaires. L'implémentation de \cite{Pegasis} se ramène à des polynômes en $X$ en exprimant les puissances paires de $Y$ comme polynômes en $X$. On approxime donc la complexité en $O(\log q)$, d'où le résultat.
\end{proof}

\subsection{Relations dans le groupe de classes d'idéaux}\label{SECTION CalculfactoClasse}
%marche aléatoire, analyse complexité, expréiences

On explique dans cette section l'approche de Broker Charles et Lauter \cite{BrokerCharlesLauter} pour trouver un idéal friable $\LR'$ équivalent à un idéal $\LR$ donné. 

\subsubsection{Famille génératrice du groupe des classes d'idéaux}

On commence par former une famille génératrice de $Cl(\OR_{\Delta})$ définie avec un système de représentants composé uniquement d'idéaux de norme vérifiant les hypothèses \ref{HypothèseInversion} et \ref{HypothèseElkies}. Pour cela, on établit la liste croissante des nombres premiers acceptables inférieurs à une borne $N$. On fixe alors une convention pour associer à chacun d'eux un idéal premier de même norme. Pour $N$ suffisamment grand on peut espérer avoir assez de représentants de classes distinctes pour obtenir une famille génératrice de $Cl(\OR_{\Delta})$.

\begin{Def}
	Soit $p$ premier décomposé dans $\OR_{\Delta}$. On fixe un algorithme déterministe prenant en entrée $p$ et $\OR_{\Delta}$ qui calcule une racine de $\OR_{\Delta}$ modulo $4p$. On définit la forme quadratique première de discriminant $\Delta$ associée à $p$ comme étant $f(p,b_{p})$, où $b_{p}\in\N$ vérifie $b_{p}^2 = \Delta [4p]$. On définit l'idéal premier associé à $p$ par $\PR_{p} = I(f(p,b_{p}))$, où l'application $I$ est définie dans la proposition \ref{applicationIidéalDeForme}. En particulier $N(\PR_{p}) = p$.
\end{Def}

\begin{Rem}
	L'idéal premier associé à $p$ est donc l'un des deux idéaux premiers au-dessus de $p$ dans $\OR_{\Delta}$. $I((p,-b_{p}))$ est le second, et représente la classe de $(\PR_{p})^{-1}$ d'après la proposition \ref{NormesId}. 
\end{Rem}

\begin{Def}
	Soit $N\in\N$. On note $P_{N} = \left\lbrace \PR_{p} : p \; premier, 0\leq p\leq N, \left( \frac{\Delta}{p}\right) = 1\right\rbrace$ la liste des idéaux premiers associés aux nombres premiers décomposés plus petits que $N$. 
\end{Def}

La proposition suivante permet de choisir la borne $N$.

\begin{Prop}[\cite{Seysen}(Th 3.3)]\label{BorneThéorique}
	
	Sous l'hypothèse de Riemann généralisée, il existe une constante $C$ telle que $Cl(\OR_{\Delta})$ soit engendré par les classes des éléments de $P_{N}$, avec 
	\begin{equation*}
		N = C\left( \log |\Delta| \right) ^{2}.
	\end{equation*}
	
\end{Prop}

\begin{Rem}
	Pour estimer la constante $C$, Seysen \cite{Seysen} cite deux résultats proposant $C= 280$ ou même $C = 2$. Cohen \cite{Coh00} propose $C = 6$. Cependant, d'après Cohen, en pratique cette borne $N$ est pessimiste et l'on conjecture une borne en $O\left( (\log |\Delta|)^{1+\epsilon}\right) $ pour tout $\epsilon > 0$. On choisira $C=2$ et on supposera $p$ suffisamment grand pour que les nombres premiers inférieurs à $N$ vérifient l'hypothèse \ref{HypothèseElkies}.
\end{Rem}

\begin{Rem}
	Comme les nombres premiers qui contredisent l'hypothèse \ref{HypothèseInversion} sont en nombre fini, les exclure de $P_{N}$ ne change pas le comportement asymptotique de la borne $N$. Prenant en compte ceci et la remarque précédente, on suppose que $P_{N}$ reste une famille génératrice si l'on ignore les idéaux de norme ne vérifiant pas l'hypothèse \ref{HypothèseInversion}. 
\end{Rem} 

En supposant $N$ choisi selon la proposition \ref{BorneThéorique}, on énonce  l'algorithme \ref{ALGO FamilleGénératrice} qui sera utilisé pour établir une liste de représentants d'une famille supposée génératrice du groupe des classes d'idéaux. 

\begin{algorithm}[H]
	\caption{Famille génératrice du groupe de classes}\label{ALGO FamilleGénératrice}
	\begin{algorithmic}[1]
		\Require Un discriminant $\Delta$, la caractéristique $p$, le conducteur maximal $f_{m}$, le cardinal $|E(\F_{q^{n}})|$. Un entier $N > 2\log|\Delta|$.
		\Ensure Une liste $L_{premiers}$ de nombres premiers vérifiant les hypothèses \ref{HypothèseInversion} et \ref{HypothèseElkies}, la liste $L_{ideaux}$ de leur idéal premier associé respectif, la liste $L_{formes}$ de leur forme première associée respective. 
		\State $L_{premiers}$, $L_{ideaux}$, $L_{formes} \gets$ \O
		\State $n \gets 2$
		\While{$n < N$}
		\If{$\left( \frac{\Delta}{n}\right) = 1$ ET $n$ ne divise pas $A$}
		\State Mettre $n$ dans $L_{premiers}$
		\State Calculer $q_{n} = (n ,b_{n})$ la forme quadratique première associée à $n$ de discriminant $\Delta$.
		\State Mettre $q_{n}$ dans $L_{formes}$
		\State Calculer $I_{n}$ l'idéal premier associé à $n$. 
		\State Mettre $I_{n}$ dans $L_{ideaux}$
		\EndIf
		\State $n \gets nextprime(n)$
		\EndWhile
		\State\Return $L_{premiers}$, $L_{ideaux}$, $L_{formes}$
	\end{algorithmic}
\end{algorithm}







\subsubsection{Écriture d'une classe dans une famille génératrice}


Soit $\LR\subset\OR_{\Delta}$ un idéal entier. Soit $L_{premiers}$, $L_{ideaux}$ et $L_{formes}$ les listes retournées par l'algorithme \ref{ALGO FamilleGénératrice}. On note $k$ leur taille, et $L_{ideaux} = \left\lbrace \PR_{1},\ldots,\PR_{k}\right\rbrace$. Alors d'après la proposition \ref{BorneThéorique}, il existe $e_{1},\ldots,e_{k}\in\Z$ tels que
\begin{equation*}
	\CL = \CP_{1}^{e_{1}}\CP_{2}^{e_{2}}\ldots\CP_{k}^{e_{k}}.
\end{equation*}

On cherche à expliciter des exposants $e_{i}$ vérifiant cette équation. Une idée naïve consiste à chercher de tels exposants au hasard. Deux faits vont permettre de rendre cela effectif:

\begin{enumerate}
	\item On peut expliciter des exposants pour certaines classes représentées par des formes quadratiques réduites ayant de bonnes propriétés.
	\item Les marches aléatoires dans le groupe de classes vont terminer de façon uniforme sur une classe quelconque.
\end{enumerate}

On commence par citer le théorème qui nous permettra d'expliciter des exposants pour certaines classes.

\begin{The}[\cite{Seysen}(Th 3.1), \cite{Coh00}(Lem 5.5.1)]\label{factorisationClasseIdéaux}
	
	Soit $\LR = id(1,a,b)\subset I(\OR_{\Delta})$ un idéal entier propre et primitif, et $f(a,b) = F(\LR)$ sa forme quadratique binaire primitive associée (proposition \ref{applicationFformeD'unIdeal}). Si $g\in\N$ et $a = p_{1}^{e_{1}}p_{2}^{e_{2}}\ldots p_{g}^{e_{g}}$ est la décomposition en facteurs premiers de $a$, alors:
	\begin{enumerate}
		\item Chaque $p_{i}$ est décomposé dans $\OR_{\Delta}$.
		\item Si $\PR_{p_{i}} = id(1,p_{i},b_{p_{i}})$ est l'idéal premier associé à $p_{i}$, alors $b = \pm b_{p_{i}} $ mod $2a$.
		\item $\LR = \PR_{p_{1}}^{\pm e_{1}}\PR_{p_{2}}^{\pm e_{2}}\ldots\PR_{p_{g}}^{\pm e_{g}}$.
		\item L'exposant de $\PR_{p_{i}}$ est $+e_{i}$ si et seulement si $b = b_{p_i} [2a]$.
	\end{enumerate}
\end{The}

\begin{Coro}\label{classesFactorisables}
	Soit $\LR\subset I(\OR_{\Delta})$ tel que $\CL$ soit représentée par une forme quadratique réduite $f(a,b)$. Si $a$ se factorise entièrement par des éléments de $L_{premiers}$ obtenus par l'algorithme \ref{ALGO FamilleGénératrice}, alors on peut expliciter des exposants d'une écriture de $\CL$ dans la famille $L_{ideaux}$ connaissant la factorisation de $a$. 
\end{Coro}

On énonce ensuite les bonnes propriétés des marches aléatoires sur les groupes de classes. On admet la théorie des graphes expanseurs, dont on cite les résultats principaux utiles par la suite.

\begin{Def}
	Soit $S$ une partie génératrice d'un groupe $G$. Le graphe de Cayley de $G$ selon $S$ est défini ainsi:
	\begin{enumerate}
		\item Les sommets sont associés chacun à un élément de $G$.
		\item Une arête relie deux sommets associés à $g$ et $g'$ si et seulement s'il existe $s\in S$ tel que $g = sg'$
	\end{enumerate}
	$S$ est supposé stable par inverse pour que le graphe ne soit pas orienté.
\end{Def}

\begin{Prop}[\cite{Jao_2009}(Lem 2.1)]\label{grapheExpanseurProba}
	
	Soit $G = Cl(\OR_{\Delta})$ et $S = L_{ideaux}\cup (L_{ideaux})^{-1}$. On pose
	\begin{equation*}
		t = O\left( \frac{\log|\Delta|}{\log\log|\Delta|}\right) 
	\end{equation*} 
	et on considère une marche aléatoire en $t$ étapes partant d'un sommet quelconque du graphe de Cayley de G selon S. Alors pour toute partie $S_{0}\subset G$, la marche aléatoire termine dans $S_{0}$ avec probabilité d'au moins $\frac{|S_{0}|}{2|G|}$. 
\end{Prop}

\begin{Rem}
	En pratique, dès que l'on enchaîne des marches aléatoires de taille de l'ordre de $\frac{\log|\Delta|}{\log\log|\Delta|}$, les résultats seront uniformément répartis parmi les éléments de $Cl(\OR_{\Delta})$. Puisque le groupe est abélien, la norme $l_{1}$ du vecteur des exposants est au pire le nombre de pas $t$ de la marche aléatoire. 
\end{Rem}

On choisit alors l'ensemble $S_{0}$ de la proposition \ref{grapheExpanseurProba} comme étant l'ensemble des classes dont on sait expliciter l'écriture dans $L_{ideaux}$ grâce au corollaire \ref{classesFactorisables}. On en déduit l'algorithme \ref{ALGO FactoClasse}, permettant de trouver des exposants $e_{i}$ d'une écriture d'une classe $\CL$ donnée dans une famille génératrice construite avec l'algorithme \ref{ALGO FamilleGénératrice}.  

\begin{algorithm}[H]
	\caption{Recherche d'exposants de la classe d'un idéal}\label{ALGO FactoClasse}
	\begin{algorithmic}[1]
		\Require Un entier positif $N$, une borne de pas de marche aléatoire $t$, un discriminant $\Delta$, la caractéristique $p$, le conducteur maximal $f_{m}$, le cardinal $|E(\F_{q^{n}})|$. Un idéal $\LR\subset\OR_{\Delta}$. 
		\Ensure Deux listes $L_{exposants}$, $L_{facteurs}$ décrivant une factorisation de $\CL$.
		\State $L_{ideaux}$, $L_{formes}$, $L_{premiers}\gets$ résultats de l'algorithme \ref{ALGO FamilleGénératrice} appliqué à $N$, $\Delta$, $p$, $f_{m}$, $|E(\F_{q^{n}})|$. 
		\State $classe_{test}\gets$ forme réduite de $\LR$
		\State $a\gets$ coefficient en $x^2$ de $classe_{test}$.
		\State $k\gets$ taille de $L_{ideaux}$
		\State $L_{x}\gets$ liste de taille $k$ initialisée à $0$
		\While{$a$ ne se factorise pas via uniquement des nombres premiers de $L_{premiers}$}
		\State $L_{x}\gets$ liste de taille $k$ initialisée à $0$
		\For{$0<i<t $}
		\State $Indice\gets$ indice aléatoire de $L_{x}$.
		\State $Increment\gets\pm1$
		\State $L_{x}\left[ Indice\right] \gets L_{x}\left[ Indice\right] + Increment$
		\EndFor
		\For{$0\leq i \leq k - 1$}
		\State $x_{i}\gets L_{x}[i]$
		\For{$1 \leq j \leq x_{i}$}
		\State $classe_{test}\gets classe_{test}L_{formes}[i]$
		\State $classe_{test}\gets classe_{test}$ réduit par l'algorithme \ref{ALGO reduction}
		\EndFor
		\EndFor
		\State $a\gets$ coefficient en $x^2$ de $classe_{test}$.
		\State Factoriser $a$ (avec ECM).
		\EndWhile
		\State $L_{y}\gets$ une liste de taille $k$ de chaque exposant des éléments de $L_{ideaux}$ dans une écriture de $I(classe_{test})$ selon le théorème \ref{factorisationClasseIdéaux}.
		\State $L_{exposants}\gets L_{y} - L_{x}$
		\State $L_{facteurs}\gets L_{ideaux}$
		\For{$e_i$ parcourant $L_{exposants}$ }
		\If{$e_i < 0$}
		\State $e_i\gets-e_i$
		\State $\left( L_{facteurs}\right) _{i} \gets $ idéal conjugué de $\left( L_{facteurs}\right) _{i}$
		\EndIf
		\EndFor
		\State \Return $L_{exposants}$, $L_{facteurs}$
	\end{algorithmic}
\end{algorithm}

\begin{Rem}
	L'algorithme \ref{ALGO FactoClasse} n'est pas exactement celui utilisé dans \cite{BrokerCharlesLauter} et \cite{JaoSoukharev}, la différence étant faite lors de l'implémentation de la marche aléatoire (ligne $7$ à $11$). Dans la section suivante, on analyse la complexité de l'algorithme énoncé. On comparera avec l'implémentation de Jao et Soukharev \cite{JaoSoukharev} dans la section \ref{SECTION BrokerExpériences}.
\end{Rem}


\subsubsection{Analyse de complexité}




L'analyse des algorithmes \ref{ALGO FamilleGénératrice} et \ref{ALGO FactoClasse} dépend du choix du paramètre $N$ bornant les nombres premiers utilisés.  

\begin{Prop}
	Soit $N$ un entier positif, $\Delta$ un discriminant. L'algorithme \ref{ALGO FamilleGénératrice} renvoie une famille supposée génératrice du groupe de classe de représentants de normes bornées par $N$ avec une complexité majorée par 
	\begin{equation*}
		O(N(\log N)^3).
	\end{equation*} 
\end{Prop}

\begin{proof}
	On parcourt tous les entiers $n$ plus petit que $N$, en testant s'ils sont premiers en $O((\log n)^4)$ (Algorithme de Miller-Rabin). Pour chaque nombre premier $n$ regardé, le calcul intermédiaire le plus coûteux est celui d'une racine $b_{n}$ de $\Delta$ modulo $4n$, en $O((\log n)^4)$ (Algorithme de Tonelli et Shanks \cite{Coh00}[1.5.1]). Le nombre de nombres premiers inférieurs à $N$ est approché par $N/\log(N)$, d'où une complexité globale majorée en $O(N(\log N)^3)$.
\end{proof}

Jao, Soukharev \cite{JaoSoukharev} et Cohen \cite{Coh00}[\S 5.5.1] proposent une optimisation du choix du paramètre $N$. Si $N$ est trop petit, on risque de ne pas avoir une famille génératrice avec l'algorithme \ref{ALGO FamilleGénératrice}, et il est plus difficile de trouver par hasard une classe facilement factorisable lors de l'algorithme \ref{ALGO FactoClasse}. Mais si $N$ est trop grand, l'algorithme \ref{ALGO FamilleGénératrice} prend plus de temps. Il sera aussi plus lent de vérifier la friabilité d'un entier dans la famille de nombres premiers $L_{premier}$, et on devra par la suite calculer des isogénies de degrés plus grand. 

\begin{Def}
	On définit la complexité sous-exponentielle en $\log|P|$, pour un certain paramètre $P$, par la fonction suivante. Pour $\alpha\in\left] 0,1 \right[$ et $z>0$, on pose
	\begin{equation*}
		L_P(\alpha,z) = e^{\left( z(\log|P|)^{\alpha}(\log\log|P|)^{1-\alpha}\right) }.
	\end{equation*}
	On dira qu'une procédure est sous-exponentielle en $\log|P|$ s'il existe des constantes $\alpha$, $z$ telles que sa complexité soit en $O(L_P(\alpha,z))$. 
	
\end{Def}

\begin{Def}
	Pour toute fonction $f$ et un certain paramètre $P$, on notera une complexité en $\tilde{O}(f(P))$ lorsque l'on omet des facteurs de tailles asymptotiquement négligeables devant $f(P)$.
\end{Def}

\begin{Def}
	Soit $\Delta$ un discriminant et $z>0$ fixé que l'on calculera par la suite. On note $N_{t} = 2(\log|\Delta|)^{2}$ et $N_{p} = L_\Delta(1/2,z)$. $N_{t}$ est appelée la borne théorique, et $N_{p}$ la borne pratique. 
\end{Def}

Les deux propositions suivantes vont permettre d'analyser la complexité de l'algorithme \ref{ALGO FactoClasse}. C'est ici que l'analyse de complexité diffère de celle de \cite{JaoSoukharev}. 

\begin{Prop}\label{espéranceExposant}
	On pose
	\begin{equation*}
		t = C_{t}\frac{\log |\Delta|}{\log \log |\Delta|}\;,\; C_{t} \geq 1\;.
	\end{equation*}
	 Soit $k$ la taille des listes retournées par l'algorithme \ref{ALGO FamilleGénératrice} appliqué à $N \geq N_{t}$. Soit $x_{i}$, $0\leq i \leq k-1$, les éléments de $L_{x}$ obtenus lors d'une itération de la boucle While de l'algorithme \ref{ALGO FactoClasse} appliqué à $N$ et $t$ (lignes $7$ à $11$). Alors les $|x_{i}|$ sont des variables aléatoires identiquement distribuées, d'espérance vérifiant
	 \begin{equation*}
	 	E(|x_{i}|)\leq\frac{C_{t}}{\log|\Delta|}. 
	 \end{equation*} 
\end{Prop}

\begin{proof}
	La liste d'exposants $L_{x}$ est de taille $k$. Le choix de la ligne $8$ est aléatoire et uniforme. Donc au cours de la boucle des lignes $7$ à $11$, chaque élément $x_{i}$ de $L_{x}$ est incrémenté en moyenne $\frac{t}{k}$ fois. On suppose qu'environ la moitié des nombres premiers plus petits que $N$ sont décomposés dans $\OR_{\Delta}$. On en déduit que $k \simeq (N/2log(N))$. En particulier si $N = N_{t}$,
	\begin{equation*}
		\frac{t}{k} \simeq C_{t}\frac{\log |\Delta|}{\log \log |\Delta|}\frac{2\log\log|\Delta|}{2(\log|\Delta|)^2} = \frac{C_{t}}{\log|\Delta|}.
	\end{equation*}
	Or $|x_{i}|$ est majoré par le nombre d'incrémentations, d'où le résultat. Si $N \geq N_{t}$, cela implique d'augmenter $k$ sans influencer $t$, donc cela ne fera que diminuer l'espérance et le résultat précédent reste une majoration. 
\end{proof}

\begin{Rem}\label{comportementListeExposant}
	On propose en pratique $C_{t} = 1$. Asymptotiquement selon $|\Delta|$, l'espérance des variables $|x_{i}|$ tend vers $0$. Ainsi on supposera que la suite $L_{x}$ contient en moyenne $t$ éléments non nuls, tous égaux à $\pm 1$. Dans la suite de l'algorithme \ref{ALGO FactoClasse}, la liste $L_{exposant}$ est obtenue en soustrayant aux éléments de $L_{x}$ les exposants correspondant issus de l'écriture de $classe_{test}$. Ces derniers sont supposés négligeables par Jao et Soukharev \cite{JaoSoukharev}, car la forme $classe_{test}$ est réduite (proposition \ref{tailleCoefReduit}). On supposera dans l'analyse de complexité que $L_{exposant}$ se comporte comme $L_{x}$. 
\end{Rem}

\begin{Rem}
	Pour éviter d'estimer $C_{t}$, on peut choisir $t = \log|\Delta|$, l'espérance des variables $|x_i|$ sera alors majorée par $\frac{\log\log|\Delta|}{\log|\Delta|}$ et tend toujours vers $0$ lorsque $|\Delta|$ augmente. 
\end{Rem}

\begin{Prop}\label{minorationNombreEssaisFacto}
	Le nombre d'itérations de la boucle While de l'algorithme \ref{ALGO FactoClasse} est minorée par 
	\begin{equation*}
		O\left( \exp \left( \left( \log\log|\Delta|  - \log\log N \right)\frac{\log|\Delta|}{2\log N} \right) \right).
	\end{equation*}
\end{Prop}

\begin{proof}
	Il s'agit d'estimer la probabilité qu'un entier $a \leq \sqrt{\frac{|\Delta|}{3}}$ soit $L_{premier}$-friable. D'après le théorème \ref{factorisationClasseIdéaux}, le coefficient $a$ ne se factorise qu'avec des nombres premiers décomposés. On s'intéresse donc à la probabilité que $a$ soit $N$-friable. Le théorème de Canfield Erdös Pomerance  \cite{Coh00}(Th 10.2.1) permet d'estimer cette probabilité comme étant 
	\begin{equation*}
		O\left( \exp \left( -\left( \log\log|\Delta|  - \log\log N  \right)\frac{\log|\Delta|}{2\log N} \right) \right).
	\end{equation*}
\end{proof}

\begin{Prop}\label{Complex iterationBoucleWhile}
	Le nombre d'opérations de chaque itération de la boucle While de l'algorithme \ref{ALGO FactoClasse} est en moyenne 
	\begin{equation*}
		O\left(\log\log|\Delta|L_{N}\left( \frac{1}{2},\frac{\sqrt{2}}{2}\right) \right).
	\end{equation*}
\end{Prop}

\begin{proof}
	Soit $t$ le nombre de pas de la marche aléatoire de l'algorithme \ref{ALGO FactoClasse}, choisi selon la proposition \ref{espéranceExposant}. Soit $k$ la taille des listes retournées par l'algorithme \ref{ALGO FamilleGénératrice}. D'après la remarque \ref{comportementListeExposant}, le calcul de $classe_{test}$ lors des lignes $13$ à $18$ demande $t$ produits, chacun entre une forme réduite et une forme première associée à un nombre premier $p_{i} \leq N$, et à chaque fois suivi d'une réduction par l'algorithme \ref{ALGO reduction}. En majorant les coefficients d'une forme réduite $(a,b,c)$ selon la proposition \ref{tailleCoefReduit},
	la complexité d'un produit de Dirichlet d'une forme réduite avec une forme première associée à $p_{i}$ est en $O(max(\log p_{i},\log|\Delta|))$ (d'après \cite{JaoSoukharev}). On note la forme obtenue $(a',b',c')$, avec $a' = ap_{i}$ et $b'\leq 2aa'$. On majore de plus $p_{i}$ par $N$. La complexité de la réduction de $(a',b',c')$ est donc majorée d'après le corollaire \ref{COMPLEX reduction} par 
	\begin{equation*}
		 O((\log N)^2(\log|\Delta|)^2).
	\end{equation*}
	Donc la complexité du calcul du produit est négligeable devant celle de la réduction qui le suit. Ainsi le calcul de $classe_{test}$ des lignes $12$ à $18$ a une complexité en
	\begin{equation*}
		O(t(\log N)^2(\log|\Delta|)^2 = O\left( \frac{(\log N)^2(\log|\Delta|)^3}{\log\log|\Delta|}\right).
	\end{equation*}
	 On note maintenant $classe_{test} = (a,b,c)$. Il reste à déterminer si $a$ est $L_{premier}$-friable. Si l'on majore les facteurs premiers de $a$ par $N$, ECM trouve un facteur de $a$ en $L_{N}(1/2,\sqrt{2}/2)$ opérations. Comme $a \leq \sqrt{|\Delta|}/3$, le nombre de diviseurs premiers distincts de $a$ est en $O(\log\log|\Delta|)$ (d'après \cite{Riesel1985}[5.8]). La factorisation de $a$ est donc de complexité
	\begin{equation*}
		O\left( \log\log|\Delta|L_{N}\left( \frac{1}{2},\frac{\sqrt{2}}{2}\right) \right), 
	\end{equation*}
	et domine asymptotiquement les calculs précédents. 
\end{proof}

On déduit de ce qui précède la complexité de l'algorithme \ref{ALGO FactoClasse}. On énonce deux propositions selon si l'on choisit $N_t$ ou $N_p$, pour prouver l'avantage asymptotique de l'optimisation proposée par \cite{JaoSoukharev} en choisissant $N_p$. 

\begin{Prop}\label{COMPLEX factoClasseTh}
	 L'algorithme \ref{ALGO FactoClasse} appliqué à $N = N_t$ admet une complexité moyenne exponentielle en $\log|\Delta|$, majorée par
	\begin{equation*}
		\tilde{O}\left( (\log l)^2 + \sqrt{|\Delta|}\right).
	\end{equation*}
	
\end{Prop}

\begin{proof}
	La seule étape de l'algorithme dépendante de $l$ est la réduction initiale de l'idéal $\LR$, en $O((\log l)^2)$. Avec $N = N_{t}$, le nombre d'itérations de la boucle While est minoré d'après la proposition \ref{minorationNombreEssaisFacto} par
	\begin{equation*}
		O\left( \exp \left(  \log\log|\Delta|   \frac{\log|\Delta|}{4\log\log|\Delta|} \right) \right) = O\left( \sqrt[4]{|\Delta|}\right).
	\end{equation*}
	La complexité est donc au moins exponentielle en $\log|\Delta|$. On peut majorer le nombre d'itérations en considérant le cas où l'on arrête la boucle While uniquement lorsque $classe_{test}$ est la classe neutre du groupe $Cl(\OR_{\Delta})$. Comme le cardinal du groupe de classe est en $O(\sqrt{|\Delta|})$ d'après la proposition \ref{cardinalGroupeClasse}, le nombre d'itérations est en moyenne $O(\sqrt{|\Delta|})$, donc exponentiel en $\log|\Delta|$. 
	
\end{proof}


\begin{Prop}\label{COMPLEX factoClasseJao}
	L'algorithme \ref{ALGO FactoClasse} appliqué à $N = N_p$ admet une complexité moyenne sous-exponentielle en $\log|\Delta|$. Plus précisément si $N_{p} = L\left(\frac{1}{2}, z \right)$ alors l'algorithme \ref{ALGO FactoClasse} admet une complexité moyenne en
	
	\begin{equation*}
		\tilde{O}\left((\log l)^2 + L_\Delta\left(\frac{1}{2}, \frac{1}{4z} \right)\right).
	\end{equation*}
	
\end{Prop}

\begin{proof}
	La seule étape de l'algorithme dépendant de $l$ est la réduction initiale de l'idéal $\LR$, en $O((\log l)^2)$. La complexité des itérations de la boucle While est d'après la proposition \ref{Complex iterationBoucleWhile} 
	\begin{equation*}
		O\left( \log\log|\Delta|L_{N_p}\left( \frac{1}{2},\frac{\sqrt{2}}{2}\right) \right) = \tilde{O}\left(  L_\Delta\left( \frac{1}{4}, \frac{\sqrt{2z}}{4} \right) \right).
	\end{equation*}
	Le nombre d'itérations est quant à lui d'après \ref{minorationNombreEssaisFacto} de l'ordre de 
	\begin{align*}
		&O\left( \exp \left( \left( \log\log|\Delta| - \log\log N_{p}  \right)\frac{\log|\Delta|}{2\log N_{p}} \right) \right) \\&= O\left( \exp \left( \log\log|\Delta|\frac{\log|\Delta|}{2\log N_{p}} \right) \right) = O\left( L_\Delta\left( \frac{1}{2},\frac{1}{4z}\right) \right).
	\end{align*}
	La complexité selon $\log|\Delta|$ de l'algorithme \ref{ALGO FactoClasse} est donc en 
	\begin{equation*}
		\tilde{O}\left(  L_\Delta\left( \frac{1}{4},\frac{\sqrt{2z}}{4}\right)L_\Delta\left( \frac{1}{2},\frac{1}{4z}\right)\right)  = \tilde{O}\left( L_\Delta\left(\frac{1}{2}, \frac{1}{4z} \right) \right).
	\end{equation*}
	
\end{proof}

\begin{Rem}\label{choixOptimalN}
	Jao et Soukharev \cite{JaoSoukharev} proposent de choisir $N = N_p$. Cependant, d'après Cohen \cite{Coh00}[5.5.1], même si la borne théorique est asymptotiquement plus petite que la borne pratique, pour des discriminants $\Delta < 10^{100}$ c'est la borne sous-exponentielle qui sera plus petite. $N_p$ n'est alors pas suffisamment grand a priori pour que l'algorithme \ref{ALGO FamilleGénératrice} appliqué à $N = N_p$ retourne une famille génératrice. De plus, lorsque l'on choisit $N = N_t$, la complexité exponentielle de l'algorithme \ref{ALGO FactoClasse} est due au fait que $N_t$ est trop petit, ce qui implique une probabilité de réussite trop faible lors des marches aléatoires. Le choix de $N$ optimal est donc
	\begin{equation*}
		N = \max(N_p,N_t).
	\end{equation*}
	On a alors asymptotiquement $N = N_p$ et la proposition \ref{COMPLEX factoClasseJao} s'applique. 
\end{Rem}


\subsection{Application au calcul d'isogénies horizontales}\label{SECTION Broker}
%Enoncés de l'algo, analyse de complexité



Soit $E\in Ell_{\F_{q}}(\OR_{\Delta})$, un nombre $l\in\N$ premier vérifiant l'hypothèse \ref{HypothèseInversion}, et un idéal $\LR$ de norme $l$. Soit $n\in\N$ et $P\in\F_{q^{n}}$. L'algorithme de Broker, Charles et Lauter \cite{BrokerCharlesLauter} calcule la $l$-isogénie de noyau $E[\LR]$, en appliquant l'algorithme \ref{ALGO FactoClasse} avec un paramètre $N$, puis en appliquant l'algorithme \ref{ALGO IsogenieEquivalenceFriable} à l'idéal équivalent friable obtenu. Le choix du paramètre $N$ sera fait selon la remarque \ref{choixOptimalN} et non pas comme proposé par Jao et Soukharev \cite{JaoSoukharev}. Dans cette section, on analyse la complexité de l'algorithme \ref{ALGO IsogenieEquivalenceFriable} lorsqu'il est appliqué à un idéal équivalent friable obtenu par l'algorithme \ref{ALGO FactoClasse}. Celle-ci est principalement dominée par le calcul et de la manipulation des polynômes modulaires. 


\begin{Prop}\label{COMPLEX equivalenceFriablePostFactoClass}
	On considère que l'entrée $\LR'$ de l'algorithme \ref{ALGO IsogenieEquivalenceFriable} est issue de l'algorithme \ref{ALGO FactoClasse} appliqué avec des paramètres $N$ et $t$ quelconques. Alors la complexité de  l'algorithme \ref{ALGO IsogenieEquivalenceFriable} est en 
	\begin{equation*}
		O\left( t\left(N^3(\log N)^{2} + N^2\log q \right) + tN^2n\log q + \left( \log l + t\log N\right)^{2}\right). 
	\end{equation*}
	
\end{Prop}

\begin{proof}
	 La liste $L_{exposant}$ retournée par l'algorithme \ref{ALGO FactoClasse} contient en moyenne $t$ exposants non nuls, tous égaux à $1$ (voir proposition \ref{espéranceExposant} et les remarques qui la suivent). Pour tout indice $i$ tel que $e_{i}$ soit non nul, il faut calculer et manipuler le polynôme modulaire $\Psi_{p_{i}}$. Le coût du calcul de $\Psi_{p_{i}}$ est, d'après Jao et Soukharev \cite{JaoSoukharev},
	\begin{equation*}
		O\left( p_{i}^{3 + \epsilon}\left(\log p_{i}\right)^{3 + \epsilon} \right),\;\forall\epsilon >0.
	\end{equation*}
	Le coût des manipulations de $\Psi_{p_{i}}$ (dérivations partielles, évaluations) est estimé par \cite{JaoSoukharev} en 
	\begin{equation*}
		O\left(p_{i}^{2 + \epsilon}\left( \log q \right)^{1 + \epsilon} \right)\;,\;\forall\epsilon >0.
	\end{equation*}
	En majorant $p_{i}$ par $N$, la complexité totale du calcul de $L_{composantes}$ est donc
	\begin{equation*}
		O\left( t \left( N^3(\log N)^{2} + N^2\log q \right) \right). 
	\end{equation*}
	Ensuite les évaluation de chaque isogénie nécessitent de mener des calculs dans le corps $\F_{q^{n}}$, car $P\in E(\F_{q^{n}})$. L'évaluation d'une isogénie de degré $p_{i}$ obtenue avec les formules de Vélu et l'algorithme BMSS nécessite $O(p_{i}^{2})$ opérations dans $\F_{q^{n}}$ (d'après \cite{JaoSoukharev}), chacune de complexité en $O(n\log q)$. On évalue au total en moyenne $t$ isogénies. En majorant $p_{i}$ par $N$, on obtient une complexité en 
	\begin{equation*}
		O(tN^2n\log q). 
	\end{equation*}
	Enfin la norme de l'idéal $\LR_{principal}$ calculée lors de la boucle For des lignes $5$ à $15$ de l'algorithme \ref{ALGO IsogenieEquivalenceFriable} est en $O(lN^t)$. Ainsi, le calcul de $\LR_{principal}$ est en $O\left( \left( \log l + t\log N\right) \right)$. De plus, la complexité de la recherche du générateur de $\LR_{principal}$ est en 
	\begin{equation*}
		O\left( \left( \log l + t\log N\right)^{2} \right). 
	\end{equation*}
\end{proof}


\begin{Coro}
	On considère que l'entrée $\LR'$ de l'algorithme \ref{ALGO IsogenieEquivalenceFriable} est issue de l'algorithme \ref{ALGO FactoClasse} appliqué avec des paramètres $N$ et $t$ choisis respectivement selon la remarque \ref{choixOptimalN} et la proposition \ref{espéranceExposant}. Alors la complexité de l'algorithme \ref{ALGO IsogenieEquivalenceFriable} est en
	\begin{equation*}
		\tilde{O}\left( (\log l)^2 + L_\Delta\left(\frac{1}{2}, \frac{1}{4z}\right) + L_\Delta\left(\frac{1}{2}, 3z\right) + \left( n+1\right) \left( \log q\right)L_\Delta\left( \frac{1}{2}, 2z\right) \right).
	\end{equation*}
	Celle-ci est minimale lorsque $z = \frac{1}{2\sqrt{3}}$, et devient
	\begin{equation*}
		\tilde{O}\left( (\log l)^2 + L_\Delta\left(\frac{1}{2}, \frac{\sqrt{3}}{2}\right) + \left( n + 1\right) \left( \log q\right) L_\Delta\left( \frac{1}{2}, \frac{1}{\sqrt{3}}\right)  \right).
	\end{equation*}
\end{Coro}

\begin{proof}
	
	Asymptotiquement, $N = N_p = L_\Delta\left(\frac{1}{2}, z \right)$. On déduit la formule énoncée des propositions \ref{COMPLEX factoClasseJao} et \ref{COMPLEX equivalenceFriablePostFactoClass}. En particulier pour tout $z$, la complexité selon $\log|\Delta|$ est
	\begin{equation*}
		\tilde{O}\left( L_\Delta\left(\frac{1}{2}, \frac{1}{4z} \right) + L_\Delta\left(\frac{1}{2}, 3z \right) \right). 
	\end{equation*}
	On choisit $z > 0$ minimisant cette complexité en équilibrant ces deux termes. Il faut donc choisir $z$ tel que
	\begin{equation*}
		\frac{1}{4z} = 3z \iff z^2 = \frac{1}{12}.
	\end{equation*}
	On choisit donc $z = \frac{1}{2\sqrt{3}}$, et on en déduit le résultat énoncé. 
	
\end{proof}

\subsection{Comparaison des implémentations}
\label{SECTION BrokerExpériences}

\subsubsection{Calculs dans les groupes de classes d'idéaux}

L'analyse de complexité de \cite{JaoSoukharev}[\S 4.4], aboutit au même résultat asymptotique, avec le même paramètre $z$ optimal. Cependant, il y a plusieurs différences entre l'algorithme proposé dans \cite{JaoSoukharev} et l'algorithmes \ref{ALGO FactoClasse}, notamment sur l'implémentation de la marche aléatoire. Dans \cite{JaoSoukharev}, on choisit indépendamment et au hasard des exposants $x_1,\ldots,x_N$, en bornant à la fois le nombre maximal de variable $x_i$ non-nulles et leur valeur absolue. Cette approche conduit à des exposants plus grands sur les idéaux de petites normes, donc à beaucoup de calculs d'isogénies intermédiaires. 

Cependant, l'utilisation de l'algorithme \ref{ALGO FactoClasse} amène à obtenir plus d'exposants non nuls, donc à calculer plus de polynômes modulaires distincts. Il est possible d'aller chercher ces derniers dans des bases de données publiques, mais en pratique pour des discriminants $|\Delta|\simeq 10^{16}$, il faudrait alors avoir accès aux polynômes $\Psi_{l}$ pour $l\simeq 3000$. La base de donnée la plus grande à ma connaissance est celle de Sutherland, et contient les polynômes $\Psi_{l}$ pour $l\simeq 1000$ premier. 

On présente les résultats d'une expérience menée dans le but de tester sur deux valeurs de $|\Delta|$ laquelle des deux méthodes est la plus efficace. Le fichier TestCLO.ipynb regroupe les codes utilisés pour cette expérience.

Soit $\Delta$ un entier négatif congru à $0$ ou $1$ modulo $4$. Pour $10^{20} < l < 10^{21}$ un nombre premier décomposé aléatoire, on calcule un idéal entier $\LR$ de $\OR_{\Delta}$ de norme $l$. Choisissant $N = \max(N_p, N_t)$, on cherche à calculer des exposants d'une écriture de $\CL$ d'abord avec la méthode de \cite{JaoSoukharev} (appelée méthode 1), puis avec l'approche de ce rapport (appelée méthode 2). On répète cela $30$ fois, avec une valeur de $l$ différente à chaque fois. Le tableau \ref{tableauexp} regroupe les moyennes des sommes des exposants, des temps de calculs, du nombre de tentatives de marches aléatoires, du nombre d'exposants non nuls et des normes maximales parmi les idéaux d'exposant non nul, le tout sur deux valeurs de $\Delta$ choisies au hasard d'ordre respectif $10^{16}$ et $10^{17}$.\\


\begin{table}[h]
	\begin{center}
		\begin{tabular}{|c|cc|cc|}
			\hline
			Discriminant & \multicolumn{2}{c|}{$\Delta = -3635657473865223$} & \multicolumn{2}{c|}{$\Delta = -48221266355996583$}  \\
			\hline
			Méthode & 1 & 2 & 1 & 2 \\
			\hline
			Somme exposants & 6266 & 12,9 & 262377 & 13 \\
			\hline
			Temps & 18,6s & 4,27s & 79s & 5,0s \\
			\hline
			Tentatives & 11,8 & 12,7 & 14,2 & 16,8 \\
			\hline
			Exposants non nuls & 12,5 & 11,9 & 15,1 & 12,2 \\
			\hline
			Norme maximale & 4515 & 4646 & 5457 & 5163 \\
			\hline
		\end{tabular}
	\end{center}
	\caption{Comparaison des méthodes 1 et 2}
	\label{tableauexp}
\end{table}

La somme des exposants correspond au nombre de calculs d'isogénies nécessaires lors de l'algorithme \ref{ALGO IsogenieEquivalenceFriable}. Le tableau \ref{tableauexp} montre l'avantage de la méthode 2 sur ce point. Le désavantage théorique est le nombre d'exposants non nuls, qui correspond au nombre de polynômes modulaires distincts à calculer. Ce désavantage est asymptotique, et pour ces valeurs de $\Delta$ la méthode 2 reste avantageuse.


\subsubsection{Calculs d'isogénies}

Le fichier Broker.ipynb regroupe mon implémentation des algorithmes \ref{ALGO BMSSCheckElkies}, \ref{ALGO IsogenieEquivalenceFriable}, \ref{ALGO FamilleGénératrice} et \ref{ALGO FactoClasse}, ainsi que mon implémentation des algorithmes décrits par Jao et Soukharev \cite{JaoSoukharev} pour les calculs dans le groupe de classes. On importe l'algorithme CheckElkies issus de l'implémentation Sage de l'article \cite{Pegasis} (voir section \ref{SECTION méthodeElkies}). Enfin, j'y ai aussi implémenté un algorithme de recherche de générateur d'un groupe de torsion (voir section \ref{SECTION generateurTorsion}), pour pouvoir calculer des isogénies directement avec l'algorithme SquareVélu de Sage. 

D'abord, j'ai testé ces algorithmes sur les exemples "small" et "medium" de l'article de Broker, Charles et Lauter \cite{BrokerCharlesLauter}. En particulier, sur l'exemple "small", ils proposent une norme de départ $l = 31$. Mais après réduction de l'idéal $\LR$ de norme $l$ avec l'algorithme \ref{ALGO reduction}, on trouve directement un idéal de norme friable, et il n'y a pas besoin de mener de marche aléatoire. Pour éviter cela, on propose de tester cette exemple avec une norme de départ $l = 59$. On trouve que les algorithmes énoncés dans le rapport, ceux de Jao et Soukharev, et ceux utilisant SquareVélu donnent les mêmes résultats, et cela en choisissant un point aléatoire de la courbe de départ.    

Ensuite, j'ai testé ces algorithmes sur l'exemple "small" de Jao est Soukharev \cite{JaoSoukharev}. Cette exemple n'est pas praticable avec l'algorithme SquareVélu. On utilise la même courbe $E$, le même idéal $\LR$ et le même point de départ $P$ que ceux de l'article, et on teste les deux implémentations. On mène les calculs sur la plateforme Matrics de l'université de Picardie Jules Vernes. Alors les algorithmes de ce rapport terminent en 357 secondes, et ceux de Jao et Soukharev en 842 secondes. Les deux algorithmes renvoient la même équation pour $E'$ et le même point $Q$. 

Cependant, l'équation trouvée diffère de celle énoncée dans l'article. De plus, les deux résultats n'ont pas le même $j$-invariant. On trouve
\begin{equation*}
	E' : y^2 = x^3 + 9704185586x + 7364944475,\;
	Q = (7467511934 ; 3986089495).
\end{equation*}
Pour vérifier leurs résultats, Jao et Soukharev calculent l'isogénie duale de $\phi_\LR:E\rightarrow E'$, c'est à dire $\phi_{\bar{\LR}} : E'\rightarrow E$, et l'évalue en $Q$. L'image de $Q$ doit alors être $lP$. On utilise la même méthode pour vérifier notre résultat. On trouve alors, avec les deux implémentations, une équation de même $j$-invariant que $E$, et à isomorphisme prêt $(x,y\mapsto l^2x, l^3y)$ on vérifie que $\phi_{\bar{\LR}}(Q) = lP$. Ainsi mon résultat semble correcte, mais on n'est pas en mesure d'expliquer la différence avec l'exemple de l'article. 